\chapter{Anaplasmosis}

\section*{Definition and Overview}
Anaplasmosis, also called human granulocytic anaplasmosis, is an infectious disease caused by the bacterium \textit{Anaplasma phagocytophilum}. It is spread to people through the bite of infected ticks, mainly the black-legged (deer) tick in the northern hemisphere.

The bacteria infect certain white blood cells called granulocytes, causing fever, headache, muscle aches, and fatigue a week or two after the bite. Most people recover fully, especially with early antibiotic treatment, but severe illness can occur, particularly in older adults or people with weakened immune systems.

\section*{Epidemiology}
Anaplasmosis occurs mainly in temperate regions of the northern hemisphere where its tick vectors live, including parts of North America, Europe, and Asia. It is most common in late spring, summer, and early autumn, when ticks are most active and people spend time outdoors.

Most cases are reported in adults, and men are affected slightly more often, probably because of greater outdoor and occupational exposure. People who hike, camp, garden, or work outdoors in tick areas are at higher risk. Cases have been increasing in many regions as tick populations spread.

\section*{Aetiology and Pathophysiology}
\begin{itemize}
    \item \textbf{Cause:} infection with the bacterium \textit{Anaplasma phagocytophilum}
    \item \textbf{Transmission:} bites from infected ticks, usually the black-legged tick; transmission typically requires the tick to be attached for many hours
    \item \textbf{How the illness develops:} the bacteria enter white blood cells called neutrophils and multiply inside them
    \item \textbf{Immune response:} the body's response to the infected cells causes fever and inflammation, and in some cases affects blood counts
    \item \textbf{Low blood counts:} infection can temporarily reduce platelet, white cell, and sometimes red cell counts
    \item \textbf{Severe disease:} the mechanism is not fully understood, but a very strong immune reaction can injure organs such as the lungs, kidneys, or liver, mainly in people who are older or immunosuppressed
\end{itemize}

\section*{Clinical Features}
\begin{itemize}
    \item Fever, chills, and severe headache, usually starting about 1 to 2 weeks after a tick bite
    \item Muscle aches, joint pain, and a general feeling of being unwell
    \item Nausea, vomiting, loss of appetite, or diarrhoea
    \item Fatigue that may last for weeks
    \item A rash in some people, although this is uncommon in anaplasmosis
    \item In severe cases, confusion, breathing difficulty, bleeding, or signs of organ failure
\end{itemize}

\section*{Differential Diagnosis}
\begin{itemize}
    \item Other tick-borne infections, such as Lyme disease, babesiosis, and tick-borne encephalitis, which can be transmitted by the same ticks
    \item Rickettsial infections such as Rocky Mountain spotted fever, which more often cause a rash
    \item Influenza, COVID-19, and other viral infections causing fever and body aches
    \item Ehrlichiosis, a similar disease caused by a related bacterium that infects different white blood cells
    \item Sepsis from other bacterial infections
    \item Other causes of prolonged fever, such as malaria, hepatitis, or urinary infection, in people with relevant exposure
\end{itemize}

\section*{When to Seek Medical Care}
See a doctor if you develop fever, headache, or muscle aches within a few weeks of a tick bite or time spent in tick areas. Mention the tick exposure, because it is key to testing.

\textbf{Call 000 (or local emergency services) for:}
\begin{itemize}
    \item Confusion, difficulty breathing, or a severe headache with a stiff neck
    \item Drowsiness, fainting, or reduced responsiveness
    \item Bruising or bleeding without a clear cause
    \item Any sudden, severe symptom that causes concern
\end{itemize}

\section*{Diagnosis and Investigations}
\begin{itemize}
    \item \textbf{History:} recent tick bite or time spent in tick habitats, and symptoms suggesting infection
    \item \textbf{Blood tests:} a full blood count may show low platelet, white cell, or red cell counts, and liver enzymes are often mildly raised
    \item \textbf{Blood film:} examination under the microscope may show characteristic clusters of bacteria inside white blood cells, although these are not always visible
    \item \textbf{PCR testing:} detects the bacterium's genetic material in the blood and is the most reliable early test
    \item \textbf{Antibody tests:} blood tests for antibodies taken during the illness and again weeks later can confirm the diagnosis
    \item \textbf{Treatment without waiting:} because tests may be negative early on, doctors often start antibiotics promptly if the illness fits the picture
\end{itemize}

\section*{Management}
\begin{itemize}
    \item \textbf{Antibiotics:} doxycycline is the standard treatment and is most effective when started early; other antibiotics may be used in special circumstances such as pregnancy
    \item \textbf{Symptom relief:} rest, fluids, and simple pain relief for fever and aches
    \item \textbf{Hospital care:} for severe illness, including intravenous antibiotics, fluids, and supportive care for affected organs
    \item \textbf{Monitoring:} blood counts and liver tests are often repeated to confirm recovery
    \item \textbf{Follow-up:} review to ensure symptoms settle, since fatigue can persist for weeks
\end{itemize}

\section*{Complications}
\begin{itemize}
    \item Severe infection affecting the lungs, kidneys, liver, or nervous system
    \item Bleeding problems from very low platelet counts
    \item Confusion, seizures, or coma in severe cases
    \item Secondary infections such as pneumonia
    \item Prolonged fatigue and weakness lasting weeks after the acute illness
    \item Death in a small proportion of severe cases, mainly in older or immunocompromised people
\end{itemize}

\section*{Prognosis and Outcomes}
Most people with anaplasmosis recover fully, and improvement is often rapid once antibiotics are started. Early treatment greatly reduces the risk of severe illness. A minority of people, especially those over 60 or with weakened immune systems, develop severe disease requiring hospital care. Even after recovery, some people feel fatigued for several weeks. With prompt diagnosis and treatment, full recovery is the usual outcome, and the disease does not generally cause lasting organ damage.

\section*{Prevention and Self-Care}
\begin{itemize}
    \item Avoid tick areas in peak season where possible, and walk in the centre of paths rather than through long grass
    \item Wear long sleeves and long trousers, and tuck trousers into socks in tick areas
    \item Use insect repellent containing DEET or other recommended agents on exposed skin
    \item Check your whole body for ticks after outdoor activity, and check children and pets
    \item Remove attached ticks promptly and correctly with fine-tipped tweezers, grasping close to the skin and pulling steadily
    \item Seek medical advice if you develop fever or flu-like symptoms after a tick bite, and tell your doctor about the bite
\end{itemize}

\section*{Sources and Further Reading}
The following resources provide reliable, peer-reviewed and patient-orientated information on anaplasmosis.

\begin{itemize}
    \item Centers for Disease Control and Prevention. \textit{Information on human anaplasmosis and tick-borne disease prevention.} https://www.cdc.gov/
    \item Mayo Clinic. \textit{Guide to the symptoms, diagnosis, and treatment of anaplasmosis.} https://www.mayoclinic.org/
    \item World Health Organization. \textit{International information on tick-borne diseases.} https://www.who.int/
    \item MSD Manual. \textit{Professional and patient information on anaplasmosis.} https://www.msdmanuals.com/
    \item MedlinePlus. \textit{Health information on anaplasmosis.} https://medlineplus.gov/
\end{itemize}
